\chapter{Analysis and discussion}\label{ch:discussion}

The preceding chapters are written in the order in which the work was done, which is also the order in which its premises were found wanting. This chapter reverses the perspective. It sets out what the records establish and what they do not, gives a physical picture that is consistent with all of them without claiming more than they support, discusses what can and cannot be said about the cosmological bounds, draws the methodological lessons we would apply again, and ends with the open questions and a concrete plan for answering them.

\section{What is established and what is not}

Four statements are, in our judgement, established, each by at least two independent calculations and each with its checks declared before the runs.

The first is that in the nonrelativistic Hermite model the work done by a resonant uniform drive follows the exact linear law within the certified windows of Table~\ref{tab:c9-win}, to $0.7\,\%$ at $\vq=0.07$ and a corresponding or smaller deficit at $0.03$ and $0.05$, with no resolved frequency shift. The residual is a seed back-reaction that scales as the square of the imposed fluctuation level and not as a property of the drive. The windows are a property of the instrument, not of the plasma. This statement is about a model whose coherent speed reaches $0.47$--$0.87c$ at the window end, and it holds for the model only.

The second is that relativity alone detunes the uniform resonance, and does so on a time scale that agrees with the loss of the resonance in HHS. The exact homogeneous relativistic kinetic equation, the relativistic grid with deterministic seeds and the relativistic Hermite model with deterministic seeds give a $10\,\%$ departure of the work from the linear law at $\wpe t=392$, $392$ and $392.5$ for $\vq=0.03$, and the grid and the exact equation give $1309.5$ and $1429$ at $\vq=10^{-3}$; the nonrelativistic grid with the same noise gives $770$--$870$ at $0.03$ and no departure at $10^{-3}$ to $3000$. The coherent mass increase of Eq.~\eqref{eq:duffing} accounts for the strong-drive number at the level of an estimate, and the thermal shift of the plasma frequency accounts for the weak-drive one, as the slow-flow model of Section~\ref{sec:relativistic} shows by giving no collapse for the weak drive without it. In the PIC the departure is earlier ($273$--$311$) because it is noise-assisted, and the Newtonian PIC shows that noise alone suffices for a $10\,\%$ departure.

The third is that an independent relativistic PIC converges, in particle number from 103 to 1648 per cell, in cell size from $0.04$ to $0.01\,c/\wpe$, in time step and in shape order, to a heating of $\simeq25\,\Eth$ at $\wpe t=1000$ in the HHS measure, and that a Newtonian pusher with the same code gives $2.5$ times more. The re-digitized HHS curve is $50.1\,\Eth$, so the plateau is half of it.

The fourth is that deterministic continuum models, with a ledger closing to $10^{-8}$ in the Hermite case, give $0.056\,\Eth$ at $\wpe t=1000$ whatever the discretization, and remain below $0.7\,\Eth$ to $2000$, a factor $445$ below the PIC plateau; and that continuum models given PIC-like noise lose positivity before the heating can develop.

Five things are not established. We do not know why our PIC plateau is a half of the HHS value. We do not know what seeds the PIC heating, beyond the exclusion of long-wavelength current noise by the mirrored loading and of a $1/N_{\rm ppc}$ noise law by the plateau. We do not know whether a positive, converged continuum initial condition with fluctuation spectrum equal to the PIC's would produce the same heating; our attempts failed on positivity, which is a property of the discretization and not a physical finding. We do not know whether the heating survives in two or three dimensions, where the instabilities that HHS invoke are not confined to the axis of the drive. And we do not know the consequence for the cosmological bounds, since no calculation here had cosmological parameters. The deterministic nonrelativistic nonlinear cascade, which the nonrelativistic windows end just before, was not resolved.

\section{A physical picture}

The picture that fits the records has two parts, which should be separated because they have different status.

The first part is established. A uniform resonant drive accelerates the electrons coherently, with speed growing as $v_q\,t/2$. Because the mass increases with the speed, the plasma frequency falls and the drive, held at the cold frequency, moves out of phase with the response. The mean-field response switches off when the accumulated phase mismatch of Eq.~\eqref{eq:slowflow}, $v_q^2t^3/64$ for a cold plasma, reaches order unity, that is at a time that scales as $v_q^{-2/3}$ and is of the order of $400$ plasma periods for $\vq=0.03$. This is a statement about a single coherent degree of freedom, it requires no fluctuations and no instability, and it explains why the physical mean field collapses at the time at which HHS see it. It is in this sense that detuning switches off the coherent pump. It is not an attractive mechanism for cosmological conclusions because it leaves the electron energy where it is: the work done by the drive up to the collapse sits in the coherent quiver and remains there, in a state in which each electron carries a large coherent momentum. Relativistic detuning alone would therefore limit the \emph{coherent} transfer without thermalizing it.

The second part is a hypothesis. The coherent state is not stable against the fluctuations present in any real plasma. Once the quiver is large the plasma is subject to the oscillating-two-stream and modulational instabilities that HHS describe, whose growth requires seeds. The PIC has a broadband seed from the finite number of particles, the continuum codes have none beyond the smooth one we gave them, and what separates the two is a factor $445$ in the heating. If the heating is seeded by fluctuations near the cell scale, as the mirrored loading and the plateau suggest and cannot prove, then the energy from the coherent quiver reaches the thermal distribution through a cascade of the kind HHS propose, in which noise at high $k$ grows fastest and the nonlinear state is dominated by it. The question that this raises is not whether such a cascade exists, since both PICs show it, but how its outcome depends on the seed spectrum. In a real early-universe plasma the seed spectrum is that of thermal fluctuations of the actual plasma, with a physical particle number per Debye volume enormously larger than a PIC can afford, so the fluctuation level at any $k$ is much smaller than a PIC's. The PIC plateau being independent of particle number over a factor of 16 would then be reassuring if the plateau had a physical cause; and worrying if it only means that a floor set by the cell is being hit, which the statement ``near the cell scale'' would imply. It is the central ambiguity of the whole study.

\section{Implications for the cosmological bounds}

A cautious reading is required. HHS conclude that the bounds from resonant conversion fail because the transfer saturates at roughly the thermal energy. What we found does not contradict the possibility that the coherent transfer is cut off, and the cut-off by relativistic detuning is robust within the one-dimensional, uniform, collisionless model. What we found does not demonstrate that the energy deposited is limited to the thermal energy either, since three further statements need to be true and none is established here.

The first is that the mean-field collapse limits the \emph{total} energy transferred, which assumes that no further energy is delivered once the coherent response is gone. In the deterministic relativistic models the energy delivered up to the collapse stays in the coherent quiver, and the model is silent about what happens to it afterwards. In the PIC the heating reaches $25\,\Eth$ and stays there to $2000$ ($28\pm3.7$), so the energy sits at the level of tens of thermal energies, and the question of the late deposition rate is not answered by a run of $2000$ plasma periods when the resonance lasts, in HHS' estimate, $10^{10}$.

The second is that the inhomogeneity needed for conversion in the early universe (the Landau--Zener passage) behaves like a uniform drive in a box. Our inhomogeneous calculations are linear and relativistic in $k$, and say that the transverse transmitted fraction follows a Landau--Zener exponent corrected by a factor $1-4.5s^2$ with $s=\eta n_c/\bar k^2$ (the notation of Section~\ref{sec:c9-r3}); in the cosmological regime the dark-photon momentum is tiny and the parameter $s$ is not necessarily small, so the standard exponent may be out of its domain. We did not test this regime.

The third is the collision frequency. Equation~\eqref{eq:c9-pss} gives a steady absorbed power inversely proportional to $\nu_{\rm eff}$ at resonance, and the numbers in HHS' Eqs.~(6) and~(7) set $\nu_{ei}\propto T_e^{-3/2}$ and an energy-exchange time $\tau_{ei}\propto T_e^{3/2}$ that grow strongly with temperature. An electron temperature $10^{2}$ times the initial one reduces $\nu_{ei}$ by $10^{3}$, and in the steady state the resonant absorption grows with the inverse of $\nu_{ei}$, which means the heating would raise the deposited power; but the amplitude transient has a time $4/\nu_{ei}$ that grows by the same factor, so the steady state would never be reached. The competition between the two is a matter of the actual Spitzer law and the actual cosmic expansion time, which our constant-$\nu$ runs cannot decide. The relativistic correction to the collision operator is also not available.

The conditional statement we are prepared to make is therefore the following. In the model of one dimension, uniform prescribed drive, constant collision rate and no cosmology, the resonant coherent transfer is cut off by relativistic mass increase at a time much shorter than the resonance time of any cosmological crossing, and the energy that reaches the thermal distribution in a relativistic PIC is of the order of tens of initial thermal energies at $\wpe t=2000$. Whether this implies that the bounds are weakened by the factors that HHS state depends on the late-time deposition, the actual collision law and the dimensionality, and the calculations here do not decide it. We have not computed, and do not give, any revised bound.

\section{Methodological lessons}

Several practices proved their value and are worth stating plainly, since they shaped what could be concluded.

The first is the certified window. A Hermite run that ends in a NaN is easy to detect; a run that stays finite and produces a plausible curve for 200 time units after it should have been abandoned is not. Defining the window in advance as the smallest of four times (refinement pairs agreeing to $10\,\%$, positivity, reaching the time, ledger closure) removed the temptation to plot until the curve looked wrong. The windows of Chapter~\ref{ch:nonrel} are short of $1000$ for the stronger drives, and the relativistic Hermite horizon of F1 fell short of its declared $2000$; both facts are in the tables, not in the discussion of the figures.

The second is declaring the gates before running and not amending them. The repository contains many failures that would have been invisible otherwise: the realized phase variance of the first coherence study, the ledger gates D and D2 of the grid, the spectrum check c of the relativistic Hermite model, which was wrong because a truncated Hermite matrix has an undamped discrete spectrum, the zero-collision rerun that differed from an old record by $6\times10^{-6}$ because the restart chunk was 200 and not 240, and the PIC convergence gates G4 and G5. In each case the failure was reported, a follow-up was declared in a new commit before it was run, and the claim was changed to what the follow-up supported. The cost is a long list of failed checks. The benefit is that a reader can see which numbers were tested against what.

The third is the independent reference. Every pair of agreeing calculations in the note (the Hermite and the grid at $392$ and $392.5$; the Landau roots from the relativistic grid against the kinetic roots; the two collision operators against one another and against the exact driven oscillator with drag; the cold-fluid and kinetic transverse conversion) is a protection against a common error, and the one place at which no independent reference exists, the heating, is where we are least sure what we know. The mistake over the HHS digitization was of the same kind from the other side: a single reading of a single curve was relied on as if it were an independent reference, and it was not.

The fourth is that agreement of a resolved calculation with a theory is silent about what happens beyond the window. The bound of Chapter~\ref{ch:nonrel} says that no saturation appears before the cascade starts, and it says nothing of the cascade. The corresponding sentences are written into each section and into the abstract of the paper, and we would keep them even if they weaken the statements.

\section{Open questions and a research plan}

The questions follow from the five items above. The plan below orders them by what each would settle and by cost. None of it has been run.

The first and cheapest step is a scan of the cell-scale noise. The PIC at 412 particles per cell and $\Delta x=0.02$ is the missing combined point, and it would directly test whether the combination of the two refinements adds to the $15\,\%$ or counts the same effect twice, as we argue. A broader scan would add a controlled injection of noise at a chosen $k$ in a PIC with a quiet start: a quiet start (mirrored velocities and regularly spaced positions) with a seed of known spectrum, amplitude and cutoff, so that one can read off the heating as a function of the seed at $k\lambda_{De}$ from $0.1$ to the cell Nyquist. If the plateau persists at a seed amplitude that is physical, the heating is robust; if it scales with the amplitude at the cell scale it is a numerical floor. The scan needs a PIC only, and the cost per run is that of the ladder above (hours on a few cores).

The second step is the shape function. The HHS code uses fifth-order splines in deposition and force interpolation; our quintic test showed a $4\,\%$ change at 412 particles per cell, but at one particle number and with a different pusher, and a combined test of fifth order with 1648 particles per cell and a smaller cell would remove the main code-level candidate for the factor of two. This is a parameter of the code we already have.

The third step is the relativistic Hermite model with a stabilised closure. The model is certified to about $1530$ and the shot-noise runs fail at once for the PIC particle number. A path to the noisy regime requires a method that preserves positivity or that damps the top of the ladder without changing the lowest moments; the options tested in the nonrelativistic case (Dai, symmetric weight, square-root form) fail their nonlinear gates at $\vq=0.1$ and have not been tried in the relativistic code. We regard the matter as unsettled, with the Hermite route not closed. A closure scaling with the local field, as in the nonrelativistic code, together with a hypercollision of the type that passed the $1\mathrm{D}2\mathrm{V}$ test, is the nearest candidate.

The fourth step is two dimensions. The instabilities invoked by HHS for the transfer of energy have a direction, the drift, and the argument for one dimension is that the fastest modes are aligned with it. A two-dimensional PIC with a mixed field would test whether the heating is the same, larger or smaller, and whether the $k=0$ resonance is detuned in the same way. The relativistic grid and Hermite models have the cost structure to do one dimension in $x$ and two in $u$ only for small $N_x$.

The fifth step is cosmological deposition. The collision law needs the Spitzer frequency $\nu_{ei}(T_e)\propto T_e^{-3/2}$ in the absorption equation, a relativistic correction, an expansion rate and cooling. With these the question becomes a closed, zero-dimensional or one-dimensional model whose deposited energy can be computed as a function of $\vq$ and $\omega_{pe}/H$ and compared with the bound; it uses Eq.~\eqref{eq:c9-drude} with a temperature-dependent $\nu$ and the heating of the PIC as an energy input, and it is the calculation that turns the present study into a statement about cosmology. It needs the answer to the heating questions first, because the input is the heating.

The remaining items, none large, are: a record of the thermal relativistic shift of the plasma frequency to confirm the weak-drive collapse time; a direct comparison of the Vay and Boris pushers in the resonant case; the Gaussian line and the combined reservoir and coherence run of Chapter~\ref{ch:nonrel}; and the missing seed levels in the back-reaction study.
